% =========================================================
% The Twelve Canonical Proof Moves in Analysis
% =========================================================
\paragraph{The Twelve Canonical Proof Moves.}

A \emph{proof move} is a recurring logical or computational technique
that appears across many different proofs in analysis. Once you learn
to name these moves, reading and writing proofs becomes a matter of
recognizing which move applies --- not inventing an argument from scratch.

The moves below are not tricks. Each one encodes a genuine mathematical
idea. But once the idea is understood, the move can be deployed almost
mechanically in the right situations.

\begin{tcolorbox}[colback=gray!6, colframe=gray!40, arc=2pt,
  left=6pt, right=6pt, top=4pt, bottom=4pt,
  title={\small\textbf{The Twelve Moves at a Glance}},
  fonttitle=\small\bfseries]
\small
\renewcommand{\arraystretch}{1.25}
\begin{tabular}{clll}
\toprule
\textbf{\#} & \textbf{Move} & \textbf{Use when \ldots} & \textbf{Level} \\
\midrule
1  & Unpack a definition & starting any proof & L1 \\
2  & Triangle inequality & need to bound $d(x,z)$ via $y$ & L2 \\
3  & $\varepsilon$-splitting & triangle ineq.\ produces a sum & L2 \\
4  & Tail argument & combining two $N$'s & L3 \\
5  & Case splitting & discrete or binary structure & L2 \\
6  & Choose $N$ explicitly & $\varepsilon$-$N$ arguments & L3 \\
7  & Contradiction & uniqueness or non-existence & L4 \\
8  & Subsequence extraction & sequence misbehaves globally & L3 \\
9  & Squeeze (sandwich) & bound $d(x_n,p)$ from above & L2 \\
10 & Inherited metric / embedding & constructing a new metric & L4 \\
11 & Open cover / finite subcover & compactness arguments & L4 \\
12 & Diagonal argument & simultaneous subsequence conditions & L3/4 \\
\bottomrule
\end{tabular}
\end{tcolorbox}

\begin{remark}[Move 1: Unpack a definition]
This is always the first move. Every analysis concept --- convergence,
Cauchy, compactness, closure, diameter --- is a compressed form of
a precise statement with quantifiers and inequalities. Unpacking the
definition turns the abstract concept into something you can work with.

Without this step, you have named objects but no machinery. With it,
you have the raw material for every subsequent move.

\textbf{Template.} Write out the definition of every term in the
hypothesis and in the goal. The hypotheses become your assumptions;
the goal becomes your target inequality or quantifier statement.
\end{remark}

\begin{remark}[Moves 2 and 3: Triangle inequality and $\varepsilon$-splitting]
These two moves almost always appear together. The triangle inequality
lets you insert an intermediate point $y$ and split $d(x,z)$ into
$d(x,y) + d(y,z)$. $\varepsilon$-splitting means choosing the bounds
on each piece so they add up to $\varepsilon$.

The intermediate point $y$ is not arbitrary: it is chosen to be something
you already have control over (a limit, a fixed center, a sequence element).
The fractions $\varepsilon/2$, $\varepsilon/3$ are chosen \emph{before}
picking $N$, so that the final bound comes out to $\varepsilon$ exactly.

\textbf{Template.} When you need $d(x_m, x_n) < \varepsilon$ and you
know each term is close to a limit $p$: insert $p$ via the triangle
inequality, budget $\varepsilon/2$ for each piece, and choose $N$ so
both pieces are bounded.
\end{remark}

\begin{remark}[Moves 4 and 6: Tail argument and choosing $N$]
The $\varepsilon$-$N$ definition of convergence requires you to \emph{produce}
an $N$. This is Move~6: work backwards from the target inequality to
determine what $N$ must satisfy, then assert its existence
(typically by the Archimedean property: for any $\varepsilon > 0$
there is an integer greater than $1/\varepsilon$).

Move~4 (the tail argument) handles the situation where two separate
conditions each give their own $N_i$. The resolution: take
$N = \max(N_1, N_2, \ldots)$. For $n \geq N$, all conditions hold
simultaneously.

These two moves are the engine of every $\varepsilon$-$N$ argument.
They appear in virtually every proof about convergent sequences.
\end{remark}

\begin{remark}[Move 5: Case splitting in analysis]
Case splitting is the same structural technique described in
Section~\ref{sec:proof-structures}, now applied to the specific
context of metric spaces. The most common case split in metric space
proofs is: $x = y$ (distance is $0$) versus $x \neq y$ (distance
is $1$ in the discrete metric, or some positive value in other spaces).

Case splitting is particularly important in verification exercises:
the triangle inequality almost always requires splitting on whether
the two outer points coincide.
\end{remark}

\begin{remark}[Move 7: Contradiction in analysis]
In analysis, contradiction is used most characteristically for
\emph{uniqueness} results: suppose two limits exist; assume they are
different; derive a contradiction from the triangle inequality and the
convergence hypothesis.

The template: suppose the goal fails (e.g., two limits exist with
$d(x,y) > 0$). Choose $\varepsilon = d(x,y)/2$. Then for large enough
$n$, both $d(x_n, x) < \varepsilon$ and $d(x_n, y) < \varepsilon$.
The triangle inequality gives $d(x,y) \leq 2\varepsilon = d(x,y)$,
a contradiction.

Notice: contradiction here uses the triangle inequality (Move~2) and
the tail argument (Move~4) internally. Proof moves nest and combine.
\end{remark}

\begin{remark}[Moves 8 and 12: Subsequence extraction and diagonal argument]
Subsequence extraction is the operation of selecting indices
$n_1 < n_2 < \cdots$ so that the restricted sequence
$\{x_{n_k}\}$ has some desired property (being monotone, being
close to a given point, etc.). It is a Level~3 operation: you are
selecting from an infinite sequence, which requires genuine sequence
machinery.

The diagonal argument applies subsequence extraction infinitely many
times and diagonalizes the result: the $k$-th term of the diagonal
sequence is $x_k^{(k)}$, where $\{x_n^{(k)}\}$ is the $k$-th
extracted subsequence. The diagonal sequence satisfies all conditions
simultaneously. This is the technique behind Arzel\`a--Ascoli and
related compactness results.
\end{remark}

\begin{remark}[Move 9: The squeeze (sandwich)]
If you can show $0 \leq d(x_n, p) \leq c_n$ where $c_n \to 0$, then
$x_n \to p$ follows immediately: any $\varepsilon$ that works for
$c_n$ also works for $d(x_n, p)$.

The squeeze is particularly useful when $d(x_n, p)$ is hard to bound
directly but can be dominated by a simpler sequence whose convergence
is known.
\end{remark}

\begin{remark}[Move 10: Inherited metric / embedding]
To construct a metric on a new space $X$, it is often easiest to
find an injective function $f : X \to Y$ into a space $Y$ whose metric
is known, and define $d_X(x,z) := d_Y(f(x), f(z))$. All four metric
axioms are then inherited from $d_Y$ via $f$. The only things to
verify are injectivity of $f$ (to ensure the identity axiom) and
whatever additional properties are required.

This is a Level~4 move because it operates on the global structure
of the space: you are choosing an ambient space and an embedding,
not just manipulating inequalities.
\end{remark}

\begin{remark}[Move 11: Open cover / finite subcover]
Compactness says: every open cover of a compact set $K$ has a finite
subcover. The power of this is that it converts infinite problems into
finite ones. If you need a single bound $\delta$ that works everywhere
on $K$, you find a $\delta_x$ at each $x$ (infinitely many), then
extract a finite subcover and take the minimum of finitely many $\delta_x$.
A finite minimum exists; a minimum over infinitely many values might not.

This move appears in all compactness arguments: uniform continuity on
compact sets, compactness implies boundedness, compactness implies
completeness in complete metric spaces.
\end{remark}
